\uuid{Kect}
\exo7id{1837}
\titre{exo7 1837}
\auteur{gourio}
\organisation{exo7}
\datecreate{2001-09-01}
\isIndication{false}
\isCorrection{false}
\chapitre{Extremum, extremum lié}
\sousChapitre{Extremum, extremum lié}
\cpge{mpsi-mp2i.fonctions-deux-variables.b, pcsi.fonctions-deux-variables.b, ptsi.fonctions-deux-variables.b, mpsi-mp2i.fonctions-deux-variables.d, pcsi.fonctions-deux-variables.d, ptsi.fonctions-deux-variables.d, bcpst1.fonctions-deux-variables.c, mp-mpi.topologie-espaces-normes.e, pc.espaces-vectoriels-normes.e, psi.espaces-vectoriels-normes.e, pt.fonctions-deux-trois-variables.b}
\module{Calcul différentiel}
\niveau{L3}
\difficulte{}

\contenu{
\question{
Soit $ A\subset {\Rr}^{2}, $ on d\'{e}finit $\int(A)$ comme
l'ensemble $ \{x\in A |\exists \rho >0,B(x,\rho )\subset A\}. $ On
supposera $ A$ ferm\'{e}e born\'{e}e et $ \int(A)\neq \emptyset .$
On suppose que $ f$ est une fonction $ C^{1 }$ sur $ A$ telle que
$ f$ est constante sur $ A\setminus \mathrm{Int}(A).  $ Montrer qu'il
existe $ z\in \mathrm{Int}(A)$ tel que :
$$\frac{\partial f}{\partial x_{1}}(z)=\frac{\partial f}{\partial x_{2}}(z)=0. $$
}
}
